% !TEX root = ../tree_notes.tex
% Appendix B. Reference implementations. Chapters derive these; the code here is
% the canonical, tested version with the same names. All compile with GCC in
% -std=c++17 and need #include <bits/stdc++.h>.
\chapter{Code Templates}\label{ch:app-templates}

These are the reference versions of the routines derived in the chapters. Copy
them once, trust them forever, and spend contest time on the idea instead of the
boilerplate.

\section{Node Types}

\begin{lstlisting}[caption={Binary tree node and n-ary tree node (Chapter 1, Chapter 2)}, label={lst:app-node}]
// Binary tree (LeetCode style).
struct TreeNode {
    int val;
    TreeNode *left, *right;
    TreeNode(int x = 0) : val(x), left(nullptr), right(nullptr) {}
};

// N-ary tree (LeetCode style): any number of children.
class Node {
public:
    int val;
    vector<Node*> children;
    Node(int x = 0) : val(x) {}
};
\end{lstlisting}

\section{General Trees from an Edge List}

Most contest trees are given as $n-1$ undirected edges, not as pointers. Store them
in an adjacency list and always pass the parent down so a DFS never walks back.

\begin{lstlisting}[caption={Build adjacency and DFS with a parent guard (Chapter 1)}, label={lst:app-adj}]
int n;
vector<vector<int>> adj;

void build(const vector<pair<int,int>>& edges) {   // nodes 0..n-1 or 1..n
    adj.assign(n, {});
    for (auto [u, v] : edges) { adj[u].push_back(v); adj[v].push_back(u); }
}

vector<int> depth_, parent_, subtree;
void dfs(int v, int p) {                            // O(n) over the whole tree
    parent_[v] = p;
    subtree[v] = 1;
    for (int to : adj[v]) {
        if (to == p) continue;                      // never go back to the parent
        depth_[to] = depth_[v] + 1;
        dfs(to, v);
        subtree[v] += subtree[to];
    }
}
\end{lstlisting}

\begin{warningbox}[The parent guard is not optional]
On a general tree stored as an adjacency list, a DFS that does not skip the parent
revisits it and loops forever. The single line \code{if (to == p) continue;} is what
turns an undirected graph traversal into a rooted-tree traversal.
\end{warningbox}

\section{Height and Depth of a Binary Tree}

The recursive form is the one to memorise (\coderef{c01-height}); the iterative form
avoids the call stack when the tree can be a path of $10^5$ nodes.

\begin{lstlisting}[caption={Recursive and iterative height (Chapter 1)}, label={lst:app-height}]
int height(TreeNode* r) {                          // height(null) = -1
    if (r == nullptr) return -1;
    return 1 + max(height(r->left), height(r->right));
}

int heightIterative(TreeNode* root) {              // post-order with an explicit stack
    if (!root) return -1;
    stack<pair<TreeNode*, int>> st;                // node, child height seen so far
    unordered_map<TreeNode*, int> h;
    st.push({root, 0});
    int best = 0;
    while (!st.empty()) {
        auto [v, state] = st.top();
        if (!v->left && !v->right) { h[v] = 0; st.pop(); best = max(best, 0); continue; }
        if (state == 0) { st.top().second = 1; if (v->left) st.push({v->left, 0}); }
        else if (state == 1) { st.top().second = 2; if (v->right) st.push({v->right, 0}); }
        else {
            int hl = v->left ? h[v->left] : -1, hr = v->right ? h[v->right] : -1;
            h[v] = 1 + max(hl, hr); best = max(best, h[v]); st.pop();
        }
    }
    return best;
}
\end{lstlisting}

\section{Level-Order Traversal (BFS)}

\begin{lstlisting}[caption={Breadth-first traversal, level by level (Chapter 3)}, label={lst:app-bfs}]
vector<vector<int>> levelOrder(TreeNode* root) {
    vector<vector<int>> levels;
    if (!root) return levels;
    queue<TreeNode*> q;
    q.push(root);
    while (!q.empty()) {
        int sz = q.size();                          // the whole current level
        vector<int> cur;
        for (int i = 0; i < sz; ++i) {
            TreeNode* v = q.front(); q.pop();
            cur.push_back(v->val);
            if (v->left)  q.push(v->left);
            if (v->right) q.push(v->right);
        }
        levels.push_back(cur);
    }
    return levels;
}
\end{lstlisting}

\section{Three Binary-Tree Traversals}

\begin{lstlisting}[caption={Preorder, inorder, postorder in one recursive pass (Chapter 3)}, label={lst:app-traversals}]
void preorder (TreeNode* r) { if (!r) return; visit(r); preorder(r->left);  preorder(r->right); }
void inorder  (TreeNode* r) { if (!r) return; inorder(r->left);  visit(r); inorder(r->right); }
void postorder(TreeNode* r) { if (!r) return; postorder(r->left); postorder(r->right); visit(r); }
\end{lstlisting}

\begin{notebox}[More templates appear as their chapters are written]
Binary search trees, rotations and balancing (\chref{c05}--\chref{c06}); heaps and
\code{priority\_queue} tricks (\chref{c07}); tries (\chref{c08}); diameter, centroid and
rerooting DP (\chref{c09}, \chref{c18}, \chref{c20}); LCA and binary lifting (\chref{c11}); segment
trees, Fenwick trees and lazy propagation (\chref{c13}, \chref{c14}, \chref{c15}); heavy--light
decomposition (\chref{c16}); DSU and MST (\chref{c28} and \chref{c17}); persistent structures
(\chref{c19}).
\end{notebox}
